\documentclass[12pt]{article}
\usepackage[utf8]{inputenc}
\usepackage{amsmath,amssymb,amsthm}
\usepackage[margin=1in]{geometry}

\newtheorem{theorem}{Theorem}
\newtheorem{lemma}[theorem]{Lemma}
\theoremstyle{definition}
\newtheorem{definition}[theorem]{Definition}

\title{Problem 334: Spilling the Beans}
\author{Project Euler Solution}
\date{}

\begin{document}
\maketitle

\section{Problem Statement}

Beans in a row of bowls are redistributed by simultaneous adjacent transfers until the configuration is stable (non-decreasing). The initial configuration is derived from the input $n$ via triangular number representations. Define $T(n)$ as the total number of bean movements. Find $T(10^7)$.

\section{Mathematical Foundation}

\begin{definition}
The $k$-th triangular number is $t_k = \frac{k(k+1)}{2}$.
\end{definition}

\begin{theorem}[Prefix-Sum Displacement]
Let $(b_0, \ldots, b_{m-1})$ be the initial and $(b_0^*, \ldots, b_{m-1}^*)$ the stable configuration. With prefix sums $S_i = \sum_{j=0}^{i} b_j$ and $S_i^* = \sum_{j=0}^{i} b_j^*$:
\[
T = \sum_{i=0}^{m-1} |S_i - S_i^*|.
\]
\end{theorem}

\begin{proof}
Each bean transfer between adjacent bowls $i$ and $i+1$ changes the prefix sum $S_i$ by $\pm 1$. The process terminates when $S_i = S_i^*$ for all $i$. Since the process follows a shortest path in the prefix-sum lattice, the total movements equal $\sum_i |S_i - S_i^*|$.
\end{proof}

\begin{lemma}[Greedy Triangular Representation]
Every non-negative integer $n$ has a unique representation as a sum of distinct triangular numbers via the greedy algorithm.
\end{lemma}

\begin{proof}
The greedy algorithm terminates since each step strictly reduces $n$. Uniqueness follows from the growth rate of triangular numbers: for any $k$, $\sum_{j=1}^{k-1} t_j < t_k$ when $k$ is sufficiently large relative to the partial sums, preventing alternative representations.
\end{proof}

\begin{theorem}[Configuration Construction]
The initial bowl configuration for input $n$ is determined by the indices appearing in the greedy triangular representation of $n$.
\end{theorem}

\begin{proof}
By construction of the problem: each triangular number $t_k$ in the representation assigns beans to specific bowls, preserving the total count.
\end{proof}

\begin{lemma}[Efficient Computation]
For a single input $n$, the greedy representation uses $O(\sqrt{n})$ triangular numbers, and $T(n)$ can be computed in $O(\sqrt{n} \log \sqrt{n})$ time.
\end{lemma}

\begin{proof}
Since $t_k \sim k^2/2$, the number of terms is $O(\sqrt{n})$. Sorting the bowl configuration takes $O(\sqrt{n} \log \sqrt{n})$, and the prefix-sum computation is linear in the number of bowls.
\end{proof}

\section{Algorithm}

\begin{enumerate}
    \item Compute the greedy triangular representation of $n$.
    \item Derive the initial bowl configuration.
    \item Sort to obtain the stable configuration.
    \item Compute prefix-sum differences and sum absolute values.
\end{enumerate}

\section{Complexity Analysis}

\begin{itemize}
    \item \textbf{Time:} $O(\sqrt{N} \log \sqrt{N})$ for a single evaluation of $T(N)$.
    \item \textbf{Space:} $O(\sqrt{N})$ for the bowl configuration.
\end{itemize}

\section{Answer}

\[
\boxed{150320021261690835}
\]

\end{document}
